\section{Literature Review}

% think needs better identification of clinical need, target population, and need for new devices, as well as the particular exercises tha patients do to combat breathlessness and their effectiveness.

Breathlessness is both subjective, and multidimensional in how patients experience it, and as such, severity is generally assessed through the self-reporting of patients, rather then through physiological measurements alone \cite{lewthwaiteHowAssessBreathlessness2021, odonnellExertionalDyspnoeaCOPD2016}. Assessment of dyspnoea is done with a variety of methods: the mMRC scale describes limitation to function in daily life, while the Borg CR10 scale is often used to describe and quantify temporary breathlessness intensity during exercise \cite{lewthwaiteHowAssessBreathlessness2021}. Because the degree of breathlessness be altered by exertion, comparison of the intensity of symptoms is most useful when ratings are obtained at set points during a standardised exercise stimulus. Simultaneous physiological measurements can then be related to patient reported ratings \cite{lewthwaiteHowAssessBreathlessness2021, odonnellExertionalDyspnoeaCOPD2016}. Despite the use of self-reporting, physiological signals still contain information that is associated with the degree of perceived dyspnoea. Dörssers et al. synchornised modified-Borg ratings with respiratory-muscle surface electromyography (EMG), respiratory rate (RR), heart rate (HR), oxygen saturation (SpO\(_2\)) and activity type during walking and cycling in 28 people with COPD \cite{dorssersDyspneaRespiratoryEMG2026}. A model containing RR, HR, activity type and SpO\(_2\) gave fitted Borg values within one point for \qty{85.97}{\percent} of observations, increasing to \qty{87.86}{\percent} after the inclusion of parasternal EMG. RR, HR and EMG were significant predictors to breathlessness intensity, while SpO\(_2\) was found to not be independently significant. However, these findings are describing fitted associations, instead of externally validated and subject independent predictions. Additionally, high quality EMG was only available in \qty{40}{\percent} of exercise segments. As such, the study supports the physiological estimation of breathlessness, but as of yet, not the reliable detection of breathlessness from a wearable device.

The feasibility of taking relevant measurements from the wrist has been investigated in other studies. Guber et al. compared wrist-based pulse oximetry with a more conventional fingertip reference in healthy participants with COPD \cite{guberWristSensorPulseOximeter2019}. Across 1,138 paired SpO\(_2\) measurements, the wrist device achieved ampltidue readings within \qty{3}{\percent} of the reference values. However, measurements that were unreliable or collected during motion were excluded from analysis, and breathlessness was not assessed. As such, evidence that physiological signals can be related to breathlessness and evidence that those signals can be measured accurately at the wrist remain separate.

Existing rehabilitation technologies demonstrate that sensing can be integrated with personalised feedback and remote supervision. COPDTrainer used therapist-derived movement parameters and smartphone inertial sensing to give real time feedback on exercise performance during unsupervised COPD rehabilitation \cite{spinaCOPDTrainerSmartphonebasedMotion2013}. Carvalho et al. combined cardiorespiratory monitoring with breathing phase feedback and physiotherapist supervision in ten people with COPD, going to show the feasibility of integrating physiological sensing and task-specific guidance within pulmonary rehabilitation \cite{carvalhoCaseStudyCOPD2026}. Otto-Yáñez et al. similarly demonstrated an eight-week supervised telerahabilitation programme incorporating repeated Borg assessment, wearable monitoring, symptom reporting and predefined recovery procedures/exercises \cite{otto-yanezPulmonaryTelerehabilitationTelemonitoring2026}. However, only nine of the fifteen enrolled participants took the study to completion, and real time HR and SpO\(_2\) monitoring during exercise was provided using primarily an O\(_2\) ring, instead of a smartwatch.

Collectively, the literature supports patient reported assessment of exercise related breathlessness, the potential value of physiological and activity signals when correlated with dyspnoea, and the technical feasibility of wearable supported pulmonary rehabilitation. In combination with the clear problem of breathlessness experienced by COPD patients, the neccessity for physician-prescribed exercises and the requirement to be able to complete these exercises at a moments notice - at home, in clinical settings, during exercise - creates a compelling argument to create and investigate new portable and effective methods of delivering these exercises and assisting users through them.